\section{Introduction}
\label{sec:tpc59-introduction}

The preceding native-affine analysis exposed an upstream support question
before either of its two Gram gates can be useful.  On the very-high face,
one writes
\begin{equation}
 n=mj+h_0=pr,
 \qquad p>y,
 \label{eq:tpc59-intro-terminal-equation}
\end{equation}
where \(p\) is prime and \(r\) is the target cofactor.  If one inserts the
convenient fixed-power cutoff
\begin{equation}
 y=X^{267/400+\eta},
 \qquad \eta>0,
 \label{eq:tpc59-intro-fixed-eta}
\end{equation}
then \(n\le X^{1+o(1)}\) gives
\begin{equation}
 r\le X^{133/400-\eta+o(1)}.
 \label{eq:tpc59-intro-small-r}
\end{equation}
This is disjoint from a dyadic block
\(R\asymp J=X^{133/400+o(1)}\).  TPC--58 therefore recorded a separate
parent-to-compatible-band loss \(\lambda_{\rm reloc}\)
\citep{WangTPC58}.

That conclusion is correct on the face declared in
\eqref{eq:tpc59-intro-fixed-eta}, but it leaves a prior optimization
question: is a fixed \(\eta\) forced by the unique-prime dictionary?  The
answer is no.  TPC--45 uses only the literal inequalities
\begin{align}
 y&>\max\{|h_0|,D_+,J_+,\Delta_M,\sqrt{N_+}\},
 \label{eq:tpc59-intro-cutoff-geometry}\\
 \frac{D_+N_+}{y}&<D_-N_-.
 \label{eq:tpc59-intro-cutoff-separation}
\end{align}
At the endpoint, \(\Delta_M=O(Q)\), every other growing term in the maximum
is \(o(Q)\), and the second inequality only imposes a bounded lower
threshold on \(y\).  A sufficiently large fixed multiple of \(Q\) therefore
satisfies both conditions.  The corresponding cofactor ceiling is again a
fixed multiple of \(J\).

This cutoff retuning removes a geometric positive-power loss, but it does
not move any atom and does not prove that the physical carrier occupies the
newly compatible window.  The distinction is central:
\begin{equation}
 \boxed{
 \text{cutoff compatibility}\ne
 \text{cofactor occupancy}\ne
 \text{terminal activation}.}
 \label{eq:tpc59-three-gates}
\end{equation}
The first is a support theorem.  The second is an orthogonal projection of
the pre-terminal carrier.  The third imposes the prime equation and all
terminal masks.  They have different sharp constants and can fail
independently.

\subsection{Main results}

Our first theorem introduces the dimensionless cutoff ratio
\begin{equation}
 \omega_X:=\frac yQ
 \label{eq:tpc59-intro-omega}
\end{equation}
and proves the exact reciprocity
\begin{equation}
 R_{\rm safe}=\frac{N_-}{2y}\asymp\frac J{\omega_X},
 \qquad
 R_{\rm max}=\frac{N_+}{y}\asymp\frac J{\omega_X}.
 \label{eq:tpc59-intro-reciprocity}
\end{equation}
A dyadic block below \(R_{\rm safe}\) is automatically contained in the
very-high face whenever it terminalizes, while a block above
\(R_{\rm max}\) is empty.  Only \(O(1)\) boundary shells lie between these
two scales.

Our second result makes occupancy exact.  In the literal lift,
\begin{equation}
 [G_X^L]_{m,r,j,\gamma}=c_{m,j,\gamma}\zeta_m,
 \label{eq:tpc59-intro-row-factorization}
\end{equation}
with no \(r\)-dependence.  For a fixed cofactor set \(\cR\), let
\(N(m)\) and \(N_{\cR}(m)\) be the total and retained cofactor
multiplicities in row \(m\), and put
\begin{equation}
 B_m:=\sum_{j,\gamma}|c_{m,j,\gamma}|^2.
 \label{eq:tpc59-intro-Bm}
\end{equation}
Then the sharp uniform energy-transfer constant is
\begin{equation}
 \boxed{
 \beta_{\cR}
 =\inf_{\cE_{\rm par}(\zeta)>0}
   \frac{\cE_{\cR}(\zeta)}{\cE_{\rm par}(\zeta)}
 =\min_{m:B_mN(m)>0}\frac{N_{\cR}(m)}{N(m)}.}
 \label{eq:tpc59-intro-beta}
\end{equation}
This is a rowwise minimax law, not an asymptotic estimate.  The inherited
carrier assumptions allow \(N_{\cR}(m)=0\) on an active row, so
\(\beta_{\cR}=0\) is possible.

Our third result identifies the terminal face exactly using the largest
prime factor of \(mj+h_0\), and partitions its raw atomic energy into
dyadic cofactor blocks.  If that face has positive energy, a maximal block
retains an \(X^{-o(1)}\) fraction.  This proves
\(\lambda_{\rm blk}=0\) after high-face activation.  It does not prove
any parent-to-face energy transfer, nor does it permit a restricted
coherent norm to be returned monotonically to the undeleted output.

\subsection{Claim boundary and organization}

The cutoff theorem, projection minimax, terminal census, and dyadic
partition are exact finite-dimensional or support statements.  Their
specialization to the literal fixed-\(h_0\) packet is L1.  A positive
cofactor-occupancy ratio, a positive terminal activation ratio, weighted
singleton occupancy, either represented native-affine Gram constant,
missing-high/off-carrier reassembly, boundary/tail reconnection, and a
lower bound for the original arithmetic defect remain open.  Nothing here
crosses the classical parity boundary
\citep{FriedlanderIwaniec2010}.

\Cref{sec:tpc59-cutoff} proves cutoff--cofactor reciprocity.
\Cref{sec:tpc59-occupancy} gives the exact projection and terminal minimax
laws.  \Cref{sec:tpc59-terminal} proves the largest-prime census and dyadic
localization.  \Cref{sec:tpc59-singleton} connects the selected block to the
undeleted singleton bridge without assuming norm monotonicity.
\Cref{sec:tpc59-endpoint,sec:tpc59-claims} give the endpoint audit, kill
tests, and revised route.
